\chapter*{Appendix: The Lean~4 Formalisation}
\addcontentsline{toc}{chapter}{Appendix: The Lean~4 Formalisation}
\label{app:lean}

Every formal claim made in this book has a machine-checked counterpart in
the Lean~4 sources distributed with the project. The formalisation
depends on Mathlib4 (release \texttt{v4.15.0}). The umbrella module
\texttt{IdeaTheory} exposes the foundational layer; subsequent
\texttt{Theorems\textsubscript{i}} files introduce the additional
operators and theorems described in the corresponding chapters; the
\texttt{Extensions} subdirectory contains the modules that ground the
applied chapters of Part~II.

\section*{Foundations}

\paragraph{\texttt{IdeaTheory.Foundations}.}
Defines the type class \texttt{IdeaTheoryStructure $\alpha$} bundling the
binary operation \texttt{op}, the identity \texttt{ident}, and the three
axioms (associativity, left identity, right identity). Proves
\texttt{ident\_unique} and the basic cancellation lemma. This is the
shared dependency of every other module.

\section*{Core theorems}

\paragraph{\texttt{IdeaTheory.Theorems1}.}
The composition theorem and its derived lemmas: powers
(\texttt{npow}), chains, sub-ideals, idempotents, cancellation, and
invertibility. Approximately seventy named results.

\paragraph{\texttt{IdeaTheory.Theorems2}.}
Coalitions and aggregation: the algebra of finite lists of ideas under
\texttt{aggregate}, committees, assemblies, the silent coalition,
coalition-equivalence and delegation.

\paragraph{\texttt{IdeaTheory.Theorems3}.}
The Aufhebung structure: maps $\sigma$, $\nu$, the derived synthesis
$\eta$, the predicates \texttt{IsPreserved}, \texttt{IsNegated},
\texttt{Sedimented}, \texttt{Sublated}, \texttt{Stable}, \texttt{Hegelian},
the triad, and the equivalence relation \texttt{aufhebianEquiv}.

\paragraph{\texttt{IdeaTheory.Theorems4}.}
The resonance pairing: the bilinear extension theorem, non-degeneracy,
the antisymmetric decomposition.

\paragraph{\texttt{IdeaTheory.Theorems5}.}
The emergence cocycle: the cocycle identity, normalisation, and the
2-cohomology classification.

\paragraph{\texttt{IdeaTheory.Theorems6}.}
The meaning curvature 2-form and its closedness; cohomology-class
invariance under structural equivalence.

\paragraph{\texttt{IdeaTheory.Theorems7}.}
Conjugation: the inner-automorphism action of invertible ideas;
preservation of $\rho$, $\sigma$, $\nu$, $\omega$.

\paragraph{\texttt{IdeaTheory.Theorems8}.}
Transmission chains: the iterated composition functor, the telescoping
identity, the chain rule for $\omega$.

\paragraph{\texttt{IdeaTheory.Theorems9}.}
Structural equivalence: the isomorphism characterisation; uniqueness up
to graded automorphism.

\paragraph{\texttt{IdeaTheory.Theorems10}--\texttt{Theorems12}.}
Graded idea algebras and the algebra of novelty: the grading theorem,
the predicates \texttt{IsNovel} and \texttt{NovelEverywhere}, advanced
reassociation lemmas, and the unifying corollaries.

\section*{Extensions used by the applied chapters}

\paragraph{\texttt{IdeaTheory.Extensions.GameTheory}.}
The general game-theoretic ground over the carrier $\Ideas$: ideas as
strategies, mixed strategies as $\RR_{\geq 0}$-cone elements summing
to one, best response, Nash equilibrium, evolutionary stability, and
the closure of equilibria under composition.

\paragraph{\texttt{IdeaTheory.Extensions.MeaningGames}.}
The Wittgenstein--Nash synthesis used by Chapter~\ref{ch:meaning-games}:
meaning games whose strategies are \emph{uses} of ideas; conventions as
Nash equilibria; effective meaning as the aggregate of a convention's
profile; family resemblance as the existence of a common idempotent;
the private-language argument as the collapse of single-speaker
games to $\unit$; salient idempotents as Lewis-style focal points;
the Kripkenstein finite-underdetermination lemma; cultural drift as
iterated meaning games; and the capstone Wittgenstein--Nash existence
theorem on finite non-empty meaning sets.

\paragraph{\texttt{IdeaTheory.Extensions.Diffusion}.}
Population-level dynamics for Chapter~\ref{ch:meaning-games} and the
sociological zoo entries of Chapters~\ref{ch:coalitions},
\ref{ch:conjugation}: the transmission step, iterated dynamics, the
weak-tie lemma, and convergence on finite populations.

\paragraph{\texttt{IdeaTheory.Extensions.RecombinantGrowth}.}
The Romer/Weitzman ground for Chapter~\ref{ch:meaning-games} and the
novelty-related zoo entries of Chapter~\ref{ch:grading}: the binary
closure of an idea pool, the iterated closure, the cardinality bound,
and the Romer non-rivalry lemma.

\paragraph{\texttt{IdeaTheory.Extensions.EmergenceLevels}.}
The emergence-cohomology ground for Chapter~\ref{ch:emergence}:
valuations, their cocycles, weak (bounded) versus strong (cohomologically
non-trivial) emergence, the coboundary lemma, and a concrete two-element
example exhibiting strong emergence.

\section*{Reproducing the build}

The project ships a \texttt{lakefile.lean} pinning Mathlib at
\texttt{v4.15.0}. After installing \texttt{elan}, the standard
\texttt{lake exe cache get \&\& lake build} sequence re-checks every
result. No \texttt{sorry} or \texttt{admit} appears in any of the
Extensions modules; the older \texttt{Theorems\textsubscript{i}} files
are imported piecewise to avoid name collisions arising from the
historical evolution of the namespace.
